\clearpage
\section{Conclusion}

The argument follows one forced velocity--pressure history from its prescribed
datum \((u_0,\nu,f)\). Differentiating the complete equation generates its
velocity, pressure, and force mixed jet. For every finite pair of temporal and
spatial orders, the datum-generated construction produces a compatible
whole-terminal rectangle containing the selected rows and their adjacent rows.
Each rectangle keeps every one of those derivative norms finite through the
chosen finite endpoint, and the compatible family forms the complete
space--time intersection on the whole maximal interval. From the same jet, and
independently of that construction, the critical-height tower and the viscous
diagonal relation supply the relationship identities on the coordinates of
every finite rectangle. Integrating the complete rectangle relations keeps
the transport, pressure, viscosity,
prescribed force, and adjacent temporal response inside those estimates
as parts of the same equation.

The adjacent rows then have continuous Sobolev representatives through the
proposed terminal time.  Their compatible values determine
\(u_*=u(T_*)\in H^\infty_\sigma\), while the forced pressure formula and the
velocity recurrence determine the complete endpoint jet from \(u_*\) and
\((\partial_t^nf(T_*))_{n\ge0}\).  Local well-posedness restarts from the
single condition \(w(T_*)=u_*\) with the same prescribed absolute-time force
\(f(t)\), written as \(f(T_*+\tau)\) in relative restart time.  Uniqueness on
the overlap joins this restart to the incoming history. Every proposed finite
maximal time therefore admits a classical continuation, contradicting
maximality. Hence the maximal time is infinite.

The resulting global forced solution is smooth, obeys the exact energy
identity with force work, and carries the decay-normalized pressure
\[
 p=R_iR_j(u_i u_j)-(-\Delta)^{-1}\nabla\cdot f.
\]
Its finite rectangles also quantify the critical-height rows, frequency
tails, Lipschitz action, vorticity amplification and dissipation, material
flow, and same-force relative-energy stability.  The small-critical theory is
an explicit quantitative refinement of this general construction: its
scale-invariant velocity stock and projected-force stock give closed formulas
for the corresponding finite rectangle bounds.

Finally, relative energy identifies every forced Leray--Hopf solution with the
same initial datum, viscosity, and prescribed force with this smooth velocity;
the common force work cancels in the comparison.  The pressure is unique up to
the unavoidable time-dependent gauge unless the decay normalization is
imposed.  Every step belongs to the one forced evolution generated by
\((u_0,\nu,f)\).
